\begin{saresult}{C.023}{Refuted}{Mixed-dimensional Cohen--Macaulay schemes require componentwise thresholds}
The seven-condition chain with one global dimension $r$ is false for
a projective Cohen--Macaulay scheme of dimension $r$ unless pure
dimension is required or the comparison thresholds are made
componentwise. A counterexample to conditions (1) and (6) of C.014
being equivalent under the literal mixed-dimensional hypothesis is
\[
 X=\mathbb P_k^2\amalg\operatorname{Spec}k,\qquad
 Y=\mathbb P_k^2,\qquad r=2,\quad m=1.
                                                        \tag{C.23.1}
\]

Here is the full corrected chain without a purity assumption.
Let $X$ be projective and Cohen--Macaulay over $k$, let $Y\subset X$
be closed, let $\mathcal O_X(1)$ be very ample, and let $m>0$.
Write the finite decomposition into open and closed subschemes
\[
 X=\coprod_sX_s,\quad \dim X_s=s\ \hbox{purely},\qquad
 Y_s=Y\cap X_s,\quad U_s=X_s\setminus Y_s,\quad d_s=s-m.
                                                        \tag{C.23.2}
\]
Omit the empty $X_s$, and let $\omega_s$ be the dualizing sheaf of
$X_s$. Write $\widehat X_s$ for its completion along $Y_s$.
The following seven conditions are equivalent. Negative sheaf
cohomology and negative Ext groups are understood to be zero.
\begin{enumerate}
\item $H^p(X\setminus Y,F)=0$ for every coherent $F$ on $X\setminus Y$
and every $p\ge m$.
\item $H^p(X\setminus Y,\mathcal O_X(-N)|_{X\setminus Y})=0$ for every
$p\ge m$ and all sufficiently large $N$.
\item For every coherent $F$ on $X$, the natural map
$H_Y^p(X,F)\to H^p(X,F)$ is bijective for $p>m$ and surjective for $p=m$.
\item Condition (3) holds for $F=\mathcal O_X(-N)$ for all sufficiently
large $N$.
\item For every $s$ and every coherent $F_s$ on $X_s$, the completion map
\[
 \operatorname{Ext}_{X_s}^q(F_s,\omega_s)\longrightarrow
 \operatorname{Ext}_{\widehat X_s}^q(\widehat F_s,\widehat\omega_s)
\]
is bijective for $q<d_s$ and injective for $q=d_s$.
\item For every $s>m$ and all sufficiently large $N$, the map
\[
 H^0(X_s,\mathcal O_{X_s}(N))\longrightarrow
 H^0(\widehat X_s,\mathcal O_{\widehat X_s}(N))
\]
is bijective and
$H^q(\widehat X_s,\mathcal O_{\widehat X_s}(N))=0$ for $0<q<d_s$.
For $s=m$, the same map on $H^0$ is required only to be injective
for all sufficiently large $N$. There is no condition for $s<m$.
\item For every $s$ and every finite locally free $E_s$ on $X_s$, the map
\[
 H^q(X_s,E_s)\longrightarrow
 H^q(\widehat X_s,\widehat E_s)
\]
is bijective for $q<d_s$ and injective for $q=d_s$.
\end{enumerate}
Thus C.014 is exactly the special case $X=X_r$ and $0<m<r$.
In the mixed-dimensional case the surjectivity clause of (4)
is automatic for large $N$ if no component has dimension $m$;
it cannot in general be deleted when such a component is present.
\end{saresult}
\begin{proof}
First check (C.23.1). The union of the plane $x_3=0$ and the point
$[0:0:0:1]$ in $\mathbb P_k^3$ realizes $X$ as a projective scheme
with homogeneous ideal
\[
 (x_0x_3,x_1x_3,x_2x_3)=(x_3)\cap(x_0,x_1,x_2).
\]
Each standard affine chart of the plane is an affine plane, and
the $x_3\ne0$ chart is a reduced point. Consequently $X$ is smooth
over $k$ on each of these disjoint open charts, hence
Cohen--Macaulay, and its dimension is two. The subscheme $Y$ is
open and closed. Its ideal is zero on $Y$ and the unit ideal on
the point, so every positive power is the same ideal and
$\widehat X=Y$. The complement is $\operatorname{Spec}k$;
therefore its coherent cohomological dimension is zero.
For every $N\ge0$, the completion map on sections is the projection
\[
 H^0(\mathbb P_k^2,\mathcal O(N))\oplus k
 \longrightarrow H^0(\mathbb P_k^2,\mathcal O(N)).
                                                        \tag{C.23.3}
\]
Its kernel is $k$. Since $r-m=1$, condition (6) of C.014 asks that
this map be bijective; it is not. This verifies the entire
counterexample, including the endpoint and arbitrarily large twists.
The attached exact computation checks the homogeneous monomial
presentation, Hilbert counts and the finite restriction matrices.

For (C.23.2), the structure map to $\operatorname{Spec}k$ is flat,
finitely presented and has a Cohen--Macaulay fibre. The precise
open-and-closed dimension decomposition is
\cite[Lemma 29.30.4, tag 02NM]{C-stacks}; equivalently, on affine
charts one applies Lemma 10.114.7, tag 00OV. Only finitely many
pieces occur because $X$ is noetherian and quasi-compact.
Every sheaf, cohomology group, formal completion and comparison
map splits over this finite decomposition.

C.012 with full-support sheaf $\mathcal O_X$ gives
(1)$\Longleftrightarrow$(2) without requiring purity. The
localization exact sequence gives
(1)$\Longleftrightarrow$(3) and
(2)$\Longleftrightarrow$(4), with the same endpoint argument
as in C.014. Coherent sheaves on the open complement extend to
coherent sheaves on $X$ as in C.014.
On each pure piece $X_s$, C.013 gives
\[
 H^p_{Y_s}(X_s,F_s)^\vee
 =\operatorname{Ext}^{s-p}_{\widehat X_s}
     (\widehat F_s,\widehat\omega_s)
\]
and the analogous ordinary-cohomology formula. It identifies
the dual of the supported-to-ordinary map with completion.
Algebraic duality is exact and faithful. Reindexing by
$q=s-p$ proves (3)$\Longleftrightarrow$(5) for every integer
$d_s$, including zero and negative values.

For a finite locally free $E_s$, take
$F_s=E_s^\vee\otimes\omega_s$. The semidualizing identity
$R\mathcal Hom(\omega_s,\omega_s)=\mathcal O_{X_s}$, and its
formal counterpart, give (5)$\Longrightarrow$(7).
Apply (7) to $E_s=\mathcal O_{X_s}(N)$ and use Serre vanishing.
When $d_s>0$ this is exactly the isomorphism on $H^0$ and
positive-degree vanishing written in (6); the endpoint injection
has zero source. When $d_s=0$ it is exactly the displayed
injection on $H^0$. When $d_s<0$ there is no condition. This
proves (7)$\Longrightarrow$(6).

Conversely, suppose (6). On each $X_s$ with $s\ge m$,
Serre vanishing turns (6) into the isomorphism/injection
comparison for positive twists through the endpoint $q=d_s$.
Use C.013 with $F_s=\omega_s(-N)$ and the same semidualizing
identity. Duality gives the supported-to-ordinary comparison
for these $F_s$, bijective for $p>m$ and surjective for $p=m$.
Localization gives
\[
 H^p(U_s,\omega_s(-N)|_{U_s})=0\qquad(p\ge m,\ N\gg0).
\]
The full-support test C.012 applied to $\omega_s$ proves
$\operatorname{cd}U_s<m$. For $s<m$, Grothendieck vanishing
\cite[Proposition 20.20.7, tag 02UZ]{C-stacks}
already gives $\operatorname{cd}U_s\le s<m$. The decomposition
is finite, so these assertions give (1).

Finally, if $s<m$, dimension vanishing gives
$H^m(X_s,\mathcal O_{X_s}(-N))=0$. If $s>m$, duality and
Serre vanishing identify its dual with
$H^{s-m}(X_s,\omega_s(N))=0$ for $N\gg0$. Thus the target in
(4) is zero when no $s=m$ piece occurs. To see why the exception
is essential, take $X=\mathbb P^2\amalg\mathbb P^1$, $m=1$,
and $Y=\mathbb P^2$. On the curve component the supported
cohomology is zero, whereas
$H^1(\mathbb P^1,\mathcal O(-N))$ has dimension $N-1$ for
$N\ge2$. Its endpoint map is not surjective.
\end{proof}
